% CC-iGPT v2 历史实验四档对比（diagnostic protocol，4 柱版，答辩简化）
% 编译: pdflatex ensemble_compare.tex   (或 latexmk -pdf)
% 嵌入论文: % CC-iGPT v2 历史实验四档对比（diagnostic protocol，4 柱版，答辩简化）
% 编译: pdflatex ensemble_compare.tex   (或 latexmk -pdf)
% 嵌入论文: % CC-iGPT v2 历史实验四档对比（diagnostic protocol，4 柱版，答辩简化）
% 编译: pdflatex ensemble_compare.tex   (或 latexmk -pdf)
% 嵌入论文: % CC-iGPT v2 历史实验四档对比（diagnostic protocol，4 柱版，答辩简化）
% 编译: pdflatex ensemble_compare.tex   (或 latexmk -pdf)
% 嵌入论文: \input{figures/ensemble_compare.tex}
%
% 数据源:
%   best / SWA / EMA: 200ep 训练日志末尾 in-training val (无 TTA, 无 ensemble)
%       best ep196 = 2.8312   (训练曲线最低点)
%       SWA last 31 = 2.8276  (Polyak 平均)
%       EMA decay 0.9998 = 2.8301
%   Ensemble + TTA: scripts/evaluate.py --ensemble best,swa,ema --tta_hflip 实测 (2026-05-27)
%       2.8296；旧 std 是无效的 batch aggregation 口径，已从图中移除
%
% 设计意图: 保留旧实验的数值溯源，不把 ensemble+TTA 当作正式单模型结果。
%   前三档 in-training 数字作为 ensemble 收益的视觉参照；全部是 diagnostic protocol 下的历史实验结果。
% 参考线:
%   v1 best+TTA 2.9035        — Phase B 历史实验（diagnostic protocol）
%   PixelSNAIL 2.85           — 文献参考线；与三模型/TTA 配置不是同协议
%   Sparse Transformer 2.80   — 文献参考线；CIFAR-10 正式比较等待单模型重评

\documentclass[tikz,border=8pt]{standalone}
\usepackage{tikz}
\usepackage{pgfplots}
\usepackage{amsmath,amssymb}
\pgfplotsset{compat=1.18}

\begin{document}
\begin{tikzpicture}[font=\footnotesize]
\begin{axis}[
    width=12cm, height=7cm,
    ybar,
    bar width=22pt,
    enlarge x limits=0.20,
    ylabel={\textbf{bits/dim (CIFAR-10 val)}},
    ylabel near ticks,
    symbolic x coords={best, SWA, EMA, Ensemble+TTA},
    xtick=data,
    x tick label style={font=\footnotesize},
    ymin=2.78, ymax=2.92,
    ytick={2.80,2.8276,2.8296,2.85,2.88,2.90,2.9035},
    yticklabel style={/pgf/number format/.cd, fixed, precision=4, zerofill, font=\scriptsize},
    axis line style={gray!55},
    tick style={gray!55},
    grid=major,
    grid style={dashed, gray!18, line width=0.3pt},
    nodes near coords,
    nodes near coords style={font=\scriptsize, /pgf/number format/.cd, fixed, precision=4, zerofill, anchor=south, yshift=2pt},
    clip=false,
]

% --- 三档 in-training (灰色, 视觉参照) ---
\addplot[draw=gray!55!black, fill=gray!25, bar shift=0pt] coordinates {
    (best,    2.8312)
    (SWA,     2.8276)
    (EMA,     2.8301)
};

% --- ensemble + TTA（橙色，diagnostic protocol 下的历史实验数字）---
\addplot[draw=orange!70!black, fill=orange!30!yellow!80, bar shift=0pt] coordinates {
    (Ensemble+TTA, 2.8296)
};

% --- 参考线 1: v1 历史实验（diagnostic protocol）2.9035 ---
\draw[densely dashed, gray!50, line width=0.4pt]
    ({rel axis cs:0,0}|-{axis cs:best,2.9035})
    -- ({rel axis cs:1,0}|-{axis cs:best,2.9035});
\node[font=\scriptsize, gray!55!black, anchor=west]
    at ({rel axis cs:1.005,0}|-{axis cs:best,2.9035}) {v1=2.9035};

% --- 参考线 2: PixelSNAIL 2.85 (协议不同，不作超越声明) ---
\draw[densely dashed, red!40, line width=0.4pt]
    ({rel axis cs:0,0}|-{axis cs:best,2.85})
    -- ({rel axis cs:1,0}|-{axis cs:best,2.85});
\node[font=\scriptsize, red!55!black, anchor=west]
    at ({rel axis cs:1.005,0}|-{axis cs:best,2.85}) {PixelSNAIL=2.85};

% --- 参考线 3: Sparse Transformer 2.80 ---
\draw[densely dashed, purple!40, line width=0.4pt]
    ({rel axis cs:0,0}|-{axis cs:best,2.80})
    -- ({rel axis cs:1,0}|-{axis cs:best,2.80});
\node[font=\scriptsize, purple!55!black, anchor=west]
    at ({rel axis cs:1.005,0}|-{axis cs:best,2.80}) {Sparse Trans=2.80};

\end{axis}

% --- 标题 ---
\node[font=\footnotesize\bfseries, anchor=south]
    at (current bounding box.north)
    {CC-iGPT v2 历史实验溯源（diagnostic protocol）— 三档 in-training 单档 vs Ensemble+TTA};

% --- 副标题 ---
\node[font=\scriptsize, gray!50!black, anchor=north, yshift=-2pt, text width=12cm, align=center]
    at (current bounding box.south)
    {\textit{N=32/d=448 (82.95M) $\cdot$ EMA 0.9998 $\cdot$ SWA last 31 ckpts (start ep170) $\cdot$ ensemble = per-token probability mixture $\cdot$ TTA = hflip 双 forward 取均\\
    灰色三档为 200ep 训练日志末尾 in-training val（无 TTA）；橙色为三成员、6-forward/image 历史实验结果（diagnostic protocol），不是正式单模型主表}};

\end{tikzpicture}
\end{document}

%
% 数据源:
%   best / SWA / EMA: 200ep 训练日志末尾 in-training val (无 TTA, 无 ensemble)
%       best ep196 = 2.8312   (训练曲线最低点)
%       SWA last 31 = 2.8276  (Polyak 平均)
%       EMA decay 0.9998 = 2.8301
%   Ensemble + TTA: scripts/evaluate.py --ensemble best,swa,ema --tta_hflip 实测 (2026-05-27)
%       2.8296；旧 std 是无效的 batch aggregation 口径，已从图中移除
%
% 设计意图: 保留旧实验的数值溯源，不把 ensemble+TTA 当作正式单模型结果。
%   前三档 in-training 数字作为 ensemble 收益的视觉参照；全部是 diagnostic protocol 下的历史实验结果。
% 参考线:
%   v1 best+TTA 2.9035        — Phase B 历史实验（diagnostic protocol）
%   PixelSNAIL 2.85           — 文献参考线；与三模型/TTA 配置不是同协议
%   Sparse Transformer 2.80   — 文献参考线；CIFAR-10 正式比较等待单模型重评

\documentclass[tikz,border=8pt]{standalone}
\usepackage{tikz}
\usepackage{pgfplots}
\usepackage{amsmath,amssymb}
\pgfplotsset{compat=1.18}

\begin{document}
\begin{tikzpicture}[font=\footnotesize]
\begin{axis}[
    width=12cm, height=7cm,
    ybar,
    bar width=22pt,
    enlarge x limits=0.20,
    ylabel={\textbf{bits/dim (CIFAR-10 val)}},
    ylabel near ticks,
    symbolic x coords={best, SWA, EMA, Ensemble+TTA},
    xtick=data,
    x tick label style={font=\footnotesize},
    ymin=2.78, ymax=2.92,
    ytick={2.80,2.8276,2.8296,2.85,2.88,2.90,2.9035},
    yticklabel style={/pgf/number format/.cd, fixed, precision=4, zerofill, font=\scriptsize},
    axis line style={gray!55},
    tick style={gray!55},
    grid=major,
    grid style={dashed, gray!18, line width=0.3pt},
    nodes near coords,
    nodes near coords style={font=\scriptsize, /pgf/number format/.cd, fixed, precision=4, zerofill, anchor=south, yshift=2pt},
    clip=false,
]

% --- 三档 in-training (灰色, 视觉参照) ---
\addplot[draw=gray!55!black, fill=gray!25, bar shift=0pt] coordinates {
    (best,    2.8312)
    (SWA,     2.8276)
    (EMA,     2.8301)
};

% --- ensemble + TTA（橙色，diagnostic protocol 下的历史实验数字）---
\addplot[draw=orange!70!black, fill=orange!30!yellow!80, bar shift=0pt] coordinates {
    (Ensemble+TTA, 2.8296)
};

% --- 参考线 1: v1 历史实验（diagnostic protocol）2.9035 ---
\draw[densely dashed, gray!50, line width=0.4pt]
    ({rel axis cs:0,0}|-{axis cs:best,2.9035})
    -- ({rel axis cs:1,0}|-{axis cs:best,2.9035});
\node[font=\scriptsize, gray!55!black, anchor=west]
    at ({rel axis cs:1.005,0}|-{axis cs:best,2.9035}) {v1=2.9035};

% --- 参考线 2: PixelSNAIL 2.85 (协议不同，不作超越声明) ---
\draw[densely dashed, red!40, line width=0.4pt]
    ({rel axis cs:0,0}|-{axis cs:best,2.85})
    -- ({rel axis cs:1,0}|-{axis cs:best,2.85});
\node[font=\scriptsize, red!55!black, anchor=west]
    at ({rel axis cs:1.005,0}|-{axis cs:best,2.85}) {PixelSNAIL=2.85};

% --- 参考线 3: Sparse Transformer 2.80 ---
\draw[densely dashed, purple!40, line width=0.4pt]
    ({rel axis cs:0,0}|-{axis cs:best,2.80})
    -- ({rel axis cs:1,0}|-{axis cs:best,2.80});
\node[font=\scriptsize, purple!55!black, anchor=west]
    at ({rel axis cs:1.005,0}|-{axis cs:best,2.80}) {Sparse Trans=2.80};

\end{axis}

% --- 标题 ---
\node[font=\footnotesize\bfseries, anchor=south]
    at (current bounding box.north)
    {CC-iGPT v2 历史实验溯源（diagnostic protocol）— 三档 in-training 单档 vs Ensemble+TTA};

% --- 副标题 ---
\node[font=\scriptsize, gray!50!black, anchor=north, yshift=-2pt, text width=12cm, align=center]
    at (current bounding box.south)
    {\textit{N=32/d=448 (82.95M) $\cdot$ EMA 0.9998 $\cdot$ SWA last 31 ckpts (start ep170) $\cdot$ ensemble = per-token probability mixture $\cdot$ TTA = hflip 双 forward 取均\\
    灰色三档为 200ep 训练日志末尾 in-training val（无 TTA）；橙色为三成员、6-forward/image 历史实验结果（diagnostic protocol），不是正式单模型主表}};

\end{tikzpicture}
\end{document}

%
% 数据源:
%   best / SWA / EMA: 200ep 训练日志末尾 in-training val (无 TTA, 无 ensemble)
%       best ep196 = 2.8312   (训练曲线最低点)
%       SWA last 31 = 2.8276  (Polyak 平均)
%       EMA decay 0.9998 = 2.8301
%   Ensemble + TTA: scripts/evaluate.py --ensemble best,swa,ema --tta_hflip 实测 (2026-05-27)
%       2.8296；旧 std 是无效的 batch aggregation 口径，已从图中移除
%
% 设计意图: 保留旧实验的数值溯源，不把 ensemble+TTA 当作正式单模型结果。
%   前三档 in-training 数字作为 ensemble 收益的视觉参照；全部是 diagnostic protocol 下的历史实验结果。
% 参考线:
%   v1 best+TTA 2.9035        — Phase B 历史实验（diagnostic protocol）
%   PixelSNAIL 2.85           — 文献参考线；与三模型/TTA 配置不是同协议
%   Sparse Transformer 2.80   — 文献参考线；CIFAR-10 正式比较等待单模型重评

\documentclass[tikz,border=8pt]{standalone}
\usepackage{tikz}
\usepackage{pgfplots}
\usepackage{amsmath,amssymb}
\pgfplotsset{compat=1.18}

\begin{document}
\begin{tikzpicture}[font=\footnotesize]
\begin{axis}[
    width=12cm, height=7cm,
    ybar,
    bar width=22pt,
    enlarge x limits=0.20,
    ylabel={\textbf{bits/dim (CIFAR-10 val)}},
    ylabel near ticks,
    symbolic x coords={best, SWA, EMA, Ensemble+TTA},
    xtick=data,
    x tick label style={font=\footnotesize},
    ymin=2.78, ymax=2.92,
    ytick={2.80,2.8276,2.8296,2.85,2.88,2.90,2.9035},
    yticklabel style={/pgf/number format/.cd, fixed, precision=4, zerofill, font=\scriptsize},
    axis line style={gray!55},
    tick style={gray!55},
    grid=major,
    grid style={dashed, gray!18, line width=0.3pt},
    nodes near coords,
    nodes near coords style={font=\scriptsize, /pgf/number format/.cd, fixed, precision=4, zerofill, anchor=south, yshift=2pt},
    clip=false,
]

% --- 三档 in-training (灰色, 视觉参照) ---
\addplot[draw=gray!55!black, fill=gray!25, bar shift=0pt] coordinates {
    (best,    2.8312)
    (SWA,     2.8276)
    (EMA,     2.8301)
};

% --- ensemble + TTA（橙色，diagnostic protocol 下的历史实验数字）---
\addplot[draw=orange!70!black, fill=orange!30!yellow!80, bar shift=0pt] coordinates {
    (Ensemble+TTA, 2.8296)
};

% --- 参考线 1: v1 历史实验（diagnostic protocol）2.9035 ---
\draw[densely dashed, gray!50, line width=0.4pt]
    ({rel axis cs:0,0}|-{axis cs:best,2.9035})
    -- ({rel axis cs:1,0}|-{axis cs:best,2.9035});
\node[font=\scriptsize, gray!55!black, anchor=west]
    at ({rel axis cs:1.005,0}|-{axis cs:best,2.9035}) {v1=2.9035};

% --- 参考线 2: PixelSNAIL 2.85 (协议不同，不作超越声明) ---
\draw[densely dashed, red!40, line width=0.4pt]
    ({rel axis cs:0,0}|-{axis cs:best,2.85})
    -- ({rel axis cs:1,0}|-{axis cs:best,2.85});
\node[font=\scriptsize, red!55!black, anchor=west]
    at ({rel axis cs:1.005,0}|-{axis cs:best,2.85}) {PixelSNAIL=2.85};

% --- 参考线 3: Sparse Transformer 2.80 ---
\draw[densely dashed, purple!40, line width=0.4pt]
    ({rel axis cs:0,0}|-{axis cs:best,2.80})
    -- ({rel axis cs:1,0}|-{axis cs:best,2.80});
\node[font=\scriptsize, purple!55!black, anchor=west]
    at ({rel axis cs:1.005,0}|-{axis cs:best,2.80}) {Sparse Trans=2.80};

\end{axis}

% --- 标题 ---
\node[font=\footnotesize\bfseries, anchor=south]
    at (current bounding box.north)
    {CC-iGPT v2 历史实验溯源（diagnostic protocol）— 三档 in-training 单档 vs Ensemble+TTA};

% --- 副标题 ---
\node[font=\scriptsize, gray!50!black, anchor=north, yshift=-2pt, text width=12cm, align=center]
    at (current bounding box.south)
    {\textit{N=32/d=448 (82.95M) $\cdot$ EMA 0.9998 $\cdot$ SWA last 31 ckpts (start ep170) $\cdot$ ensemble = per-token probability mixture $\cdot$ TTA = hflip 双 forward 取均\\
    灰色三档为 200ep 训练日志末尾 in-training val（无 TTA）；橙色为三成员、6-forward/image 历史实验结果（diagnostic protocol），不是正式单模型主表}};

\end{tikzpicture}
\end{document}

%
% 数据源:
%   best / SWA / EMA: 200ep 训练日志末尾 in-training val (无 TTA, 无 ensemble)
%       best ep196 = 2.8312   (训练曲线最低点)
%       SWA last 31 = 2.8276  (Polyak 平均)
%       EMA decay 0.9998 = 2.8301
%   Ensemble + TTA: scripts/evaluate.py --ensemble best,swa,ema --tta_hflip 实测 (2026-05-27)
%       2.8296；旧 std 是无效的 batch aggregation 口径，已从图中移除
%
% 设计意图: 保留旧实验的数值溯源，不把 ensemble+TTA 当作正式单模型结果。
%   前三档 in-training 数字作为 ensemble 收益的视觉参照；全部是 diagnostic protocol 下的历史实验结果。
% 参考线:
%   v1 best+TTA 2.9035        — Phase B 历史实验（diagnostic protocol）
%   PixelSNAIL 2.85           — 文献参考线；与三模型/TTA 配置不是同协议
%   Sparse Transformer 2.80   — 文献参考线；CIFAR-10 正式比较等待单模型重评

\documentclass[tikz,border=8pt]{standalone}
\usepackage{tikz}
\usepackage{pgfplots}
\usepackage{amsmath,amssymb}
\pgfplotsset{compat=1.18}

\begin{document}
\begin{tikzpicture}[font=\footnotesize]
\begin{axis}[
    width=12cm, height=7cm,
    ybar,
    bar width=22pt,
    enlarge x limits=0.20,
    ylabel={\textbf{bits/dim (CIFAR-10 val)}},
    ylabel near ticks,
    symbolic x coords={best, SWA, EMA, Ensemble+TTA},
    xtick=data,
    x tick label style={font=\footnotesize},
    ymin=2.78, ymax=2.92,
    ytick={2.80,2.8276,2.8296,2.85,2.88,2.90,2.9035},
    yticklabel style={/pgf/number format/.cd, fixed, precision=4, zerofill, font=\scriptsize},
    axis line style={gray!55},
    tick style={gray!55},
    grid=major,
    grid style={dashed, gray!18, line width=0.3pt},
    nodes near coords,
    nodes near coords style={font=\scriptsize, /pgf/number format/.cd, fixed, precision=4, zerofill, anchor=south, yshift=2pt},
    clip=false,
]

% --- 三档 in-training (灰色, 视觉参照) ---
\addplot[draw=gray!55!black, fill=gray!25, bar shift=0pt] coordinates {
    (best,    2.8312)
    (SWA,     2.8276)
    (EMA,     2.8301)
};

% --- ensemble + TTA（橙色，diagnostic protocol 下的历史实验数字）---
\addplot[draw=orange!70!black, fill=orange!30!yellow!80, bar shift=0pt] coordinates {
    (Ensemble+TTA, 2.8296)
};

% --- 参考线 1: v1 历史实验（diagnostic protocol）2.9035 ---
\draw[densely dashed, gray!50, line width=0.4pt]
    ({rel axis cs:0,0}|-{axis cs:best,2.9035})
    -- ({rel axis cs:1,0}|-{axis cs:best,2.9035});
\node[font=\scriptsize, gray!55!black, anchor=west]
    at ({rel axis cs:1.005,0}|-{axis cs:best,2.9035}) {v1=2.9035};

% --- 参考线 2: PixelSNAIL 2.85 (协议不同，不作超越声明) ---
\draw[densely dashed, red!40, line width=0.4pt]
    ({rel axis cs:0,0}|-{axis cs:best,2.85})
    -- ({rel axis cs:1,0}|-{axis cs:best,2.85});
\node[font=\scriptsize, red!55!black, anchor=west]
    at ({rel axis cs:1.005,0}|-{axis cs:best,2.85}) {PixelSNAIL=2.85};

% --- 参考线 3: Sparse Transformer 2.80 ---
\draw[densely dashed, purple!40, line width=0.4pt]
    ({rel axis cs:0,0}|-{axis cs:best,2.80})
    -- ({rel axis cs:1,0}|-{axis cs:best,2.80});
\node[font=\scriptsize, purple!55!black, anchor=west]
    at ({rel axis cs:1.005,0}|-{axis cs:best,2.80}) {Sparse Trans=2.80};

\end{axis}

% --- 标题 ---
\node[font=\footnotesize\bfseries, anchor=south]
    at (current bounding box.north)
    {CC-iGPT v2 历史实验溯源（diagnostic protocol）— 三档 in-training 单档 vs Ensemble+TTA};

% --- 副标题 ---
\node[font=\scriptsize, gray!50!black, anchor=north, yshift=-2pt, text width=12cm, align=center]
    at (current bounding box.south)
    {\textit{N=32/d=448 (82.95M) $\cdot$ EMA 0.9998 $\cdot$ SWA last 31 ckpts (start ep170) $\cdot$ ensemble = per-token probability mixture $\cdot$ TTA = hflip 双 forward 取均\\
    灰色三档为 200ep 训练日志末尾 in-training val（无 TTA）；橙色为三成员、6-forward/image 历史实验结果（diagnostic protocol），不是正式单模型主表}};

\end{tikzpicture}
\end{document}
